\documentclass[12pt]{article}
\usepackage[margin=1in]{geometry}
\usepackage{amsmath, amssymb, amsfonts, bm}
\usepackage{enumitem}
\setlist[enumerate]{nosep}
\setlength{\parskip}{0.7em}
\setlength{\parindent}{0pt}

\begin{document}

\begin{center}
{\Large \textbf{Solutions -- Final Exam}} \\[0.9em]
\textbf{ECON 320 -- Introduction to Mathematical Economics} \\[0.7em]

December 10, 2025
 \\[0.8em]
Instructor: Muhebullah Karimzada
\end{center}

\hrule

\textbf{Instructions:}
\begin{itemize}
\item You are allowed \textbf{120 minutes} to complete the exam.
    \item Show all work. Unsupported answers receive little credit.

    \item Unless stated otherwise, all variables are real and differentiable as needed.
\end{itemize}

\hrule
\section*{Multiple Choice Questions (20 points)}

Each question is worth 2 points. Correct options are indicated and briefly justified.

\textbf{1.} \(Q = \sqrt{P+10} - I\).

\[
\frac{\partial Q}{\partial I} = -1.
\]

\textbf{Answer:} (b) \(-1\).

\medskip

\textbf{2.} \(z = x^2 y^3\).

\[
z_x = 2x y^3,\qquad z_y = 3x^2 y^2,
\]
\[
dz = z_x\,dx + z_y\,dy = 2x y^3\,dx + 3x^2 y^2\,dy.
\]

\textbf{Answer:} (a).

\medskip

\textbf{3.} \(f(x)=4x - x^3\).

\[
f'(x)=4 - 3x^2 = 0 \Rightarrow x^2=\frac{4}{3} \Rightarrow x=\pm \sqrt{\frac{4}{3}}.
\]
Second derivative:
\[
f''(x) = -6x.
\]
At \(x=\sqrt{4/3}>0\), \(f''(x) < 0 \Rightarrow\) local maximum.

\textbf{Answer:} (a) local maximum.

\medskip

\textbf{4.} \(Y = A K^{0.5} L^{0.5}\).

In Cobb--Douglas, elasticity of \(Y\) w.r.t.\ \(K\) equals the exponent on \(K\), here \(0.5\).

\textbf{Answer:} (a) \(0.5\).

\medskip

\textbf{5.} \(f(x,y)=3x^2 + 4xy + y^2\).

\[
f_{xx}=6,\quad f_{yy}=2,\quad f_{xy}=f_{yx}=4,
\]
\[
\det(H)=6\cdot 2 - 4^2 = 12 - 16 = -4 < 0,
\]
so the stationary point (if any) is a saddle.

\textbf{Answer:} (a) saddle point.

\medskip

\textbf{6.} \(\max U(x,y)=xy\) s.t.\ \(x+y=10\).

The Lagrange multiplier \(\lambda\) measures the marginal increase in maximum utility if the right-hand side of the constraint (10) is increased: shadow value of relaxing the constraint.

\textbf{Answer:} (c) marginal utility of relaxing the constraint.

\medskip

\textbf{7.} \(\max_x f(x,a)=ax - x^2\).

Envelope theorem: 
\[
\frac{dv}{da} = \left.\frac{\partial f}{\partial a}\right|_{x=x^\ast(a)} = x^\ast.
\]

\textbf{Answer:} (a) \(x\).

\medskip

\textbf{8.} Demand \(P(Q)=20-Q\), price \(P=10\).

At \(P=10\), \(10=20-Q \Rightarrow Q=10\).  
Consumer surplus is area of triangle:
\[
\text{CS} = \frac{1}{2}(\text{choke price}-P)\cdot Q
= \frac{1}{2}(20-10)\cdot 10 = \frac{1}{2}\cdot 10\cdot 10 = 50.
\]

\textbf{Answer:} (b) 50.

\medskip

\textbf{9.} \(\dfrac{dy}{dt}=4y\), \(y(0)=3\).

Solve:
\[
y(t)=Ce^{4t},\quad y(0)=3 \Rightarrow C=3 \Rightarrow y(t)=3e^{4t}.
\]

\textbf{Answer:} (b) \(y=3e^{4t}\).

\medskip
\newpage
\textbf{10.} \(Y(t)=K^{0.3} L^{0.7} e^{0.02 t}\) with \(K,L\) constant.

\[
\frac{dY}{dt} = K^{0.3} L^{0.7} \cdot 0.02 e^{0.02 t} = 0.02 Y.
\]

\textbf{Answer:} (a) \(\frac{dY}{dt}=0.02Y\).

\bigskip

\section*{Problems (80 points total)}

\subsection*{Problem 1 (16 points)}

A firm produces
\[
Q = x^{1/2} y^{1/3}.
\]

\textbf{(a) Compute \(dQ\). (5 pts)}

Partial derivatives:
\[
Q_x = \frac{\partial Q}{\partial x}
= \frac{1}{2} x^{-1/2} y^{1/3},
\qquad 
Q_y = \frac{\partial Q}{\partial y}
= \frac{1}{3} x^{1/2} y^{-2/3}.
\]
Total differential:
\[
dQ = Q_x\,dx + Q_y\,dy
= \frac{1}{2} x^{-1/2} y^{1/3}\,dx
+ \frac{1}{3} x^{1/2} y^{-2/3}\,dy.
\]

\textbf{(b) If the firm controls only \(x\), compute \(dQ/dx\). (3 pts)}

Holding \(y\) fixed:
\[
\frac{dQ}{dx} = Q_x = \frac{1}{2} x^{-1/2} y^{1/3}.
\]

\textbf{(c) Profit \(\pi(x)=40Q - x - y - 5\), find optimal \(x\) (with \(y\) fixed). (5 pts)}

Substitute \(Q\):
\[
\pi(x) = 40 x^{1/2} y^{1/3} - x - y - 5.
\]
Differentiate w.r.t.\ \(x\) (treating \(y\) as constant):
\[
\pi'(x) = 40 \cdot \frac{1}{2} x^{-1/2} y^{1/3} - 1
= 20 x^{-1/2} y^{1/3} - 1.
\]
FOC:
\[
20 x^{-1/2} y^{1/3} - 1 = 0
\quad\Rightarrow\quad
20 x^{-1/2} y^{1/3} = 1.
\]
Solve for \(x\). Rearrange:
\[
x^{-1/2} = \frac{1}{20 y^{1/3}}
\quad\Rightarrow\quad
x^{1/2} = 20 y^{1/3}.
\]
Squaring:
\[
x^\ast = (20 y^{1/3})^2 = 400 y^{2/3}.
\]

\textbf{(d) Maximum or minimum? (3 pts)}

Second derivative:
\[
\pi''(x) = \frac{d}{dx}\left(20 x^{-1/2} y^{1/3} - 1\right)
= 20 y^{1/3} \cdot \left(-\frac{1}{2} x^{-3/2}\right)
= -10 y^{1/3} x^{-3/2}.
\]
For \(x>0\) and \(y>0\), \(x^{-3/2}>0\), \(y^{1/3}>0\), hence
\[
\pi''(x^\ast) < 0,
\]
so \(\pi\) is locally concave and \(x^\ast\) is a local maximum (and global maximum on \(x>0\)).

\bigskip

\subsection*{Problem 2 (16 points)}

\[
R(t)=500 e^{0.03 t}(1+t^{1/2}).
\]

\textbf{(a) Compute \(\frac{dR}{dt}\). (5 pts)}

Write \(R(t) = 500 u(t) v(t)\) with
\[
u(t) = e^{0.03 t}, \quad v(t) = 1 + t^{1/2}.
\]
Then
\[
u'(t) = 0.03 e^{0.03 t},\quad v'(t) = \frac{1}{2} t^{-1/2}.
\]
Product rule:
\[
\frac{dR}{dt} = 500 \big(u'v + uv'\big)
= 500\left[0.03 e^{0.03 t}(1+t^{1/2}) + e^{0.03 t}\frac{1}{2}t^{-1/2}\right].
\]
Factor:
\[
\frac{dR}{dt} 
= 500 e^{0.03 t}\left[0.03(1+t^{1/2}) + \frac{1}{2}t^{-1/2}\right].
\]

\textbf{(b) Growth rate \(g(t) = \frac{1}{R}\frac{dR}{dt}\). (4 pts)}

\[
R(t) = 500 e^{0.03 t}(1+t^{1/2}),
\]
so
\[
g(t) 
= \frac{\frac{dR}{dt}}{R}
= \frac{500 e^{0.03 t}\left[0.03(1+t^{1/2}) + \frac{1}{2}t^{-1/2}\right]}{500 e^{0.03 t} (1+t^{1/2})}.
\]
Cancel the common factor:
\[
g(t)
= \frac{0.03(1+t^{1/2}) + \frac{1}{2}t^{-1/2}}{1+t^{1/2}}
= 0.03 + \frac{1}{2}\, \frac{t^{-1/2}}{1+t^{1/2}}.
\]

\textbf{(c) Elasticity of \(R\) w.r.t.\ time. (4 pts)}

Elasticity:
\[
E_{R,t} = \frac{dR}{dt}\cdot \frac{t}{R} = g(t)\, t.
\]
Hence
\[
E_{R,t}(t) 
= t\left[0.03 + \frac{1}{2}\frac{t^{-1/2}}{1+t^{1/2}}\right]
= 0.03 t + \frac{t^{1/2}}{2(1+t^{1/2})}.
\]

\textbf{(d) Accelerating or decelerating growth? (3 pts)}

Observe:
\[
g(t) = 0.03 + \frac{1}{2}\frac{t^{-1/2}}{1+t^{1/2}}.
\]
The second term is positive for \(t>0\) but \emph{decreases} as \(t\) increases, and
\[
\lim_{t\to\infty} g(t) = 0.03.
\]
So the growth rate declines toward a constant \(3\%\). Revenue growth is \emph{decelerating}, approaching a steady growth rate.

\bigskip

\subsection*{Problem 3 (16 points)}

\[
f(x,y)=20x - 4x^2 + 10y - 3y^2 - xy.
\]

\textbf{(a) First-order partial derivatives. (4 pts)}

\[
f_x = 20 - 8x - y,
\qquad
f_y = 10 - 6y - x.
\]

\textbf{(b) Solve the first-order conditions. (4 pts)}

Set \(f_x=f_y=0\):
\[
20 - 8x - y = 0 \quad \Rightarrow \quad y = 20 - 8x,
\]
\[
10 - 6y - x = 0 \quad \Rightarrow \quad 6y + x = 10 \Rightarrow y = \frac{10 - x}{6}.
\]
Equate the two expressions for \(y\):
\[
20 - 8x = \frac{10 - x}{6}.
\]
Multiply both sides by 6:
\[
120 - 48x = 10 - x
\quad\Rightarrow\quad
120 - 10 = 48x - x
\quad\Rightarrow\quad
110 = 47x.
\]
Thus
\[
x^\ast = \frac{110}{47}.
\]
Then
\[
y^\ast = 20 - 8x^\ast
= 20 - 8\left(\frac{110}{47}\right)
= \frac{20\cdot 47 - 8\cdot 110}{47}
= \frac{940 - 880}{47}
= \frac{60}{47}.
\]

\textbf{(c) Hessian and determinant at the critical point. (4 pts)}

Second derivatives:
\[
f_{xx} = -8,\quad f_{yy} = -6,\quad f_{xy}=f_{yx} = -1.
\]
Hessian:
\[
H = \begin{bmatrix}
-8 & -1\\
-1 & -6
\end{bmatrix}.
\]
Determinant:
\[
\det(H) = (-8)(-6) - (-1)^2 = 48 - 1 = 47 > 0.
\]

\textbf{(d) Classification. (4 pts)}

We have \(f_{xx}=-8<0\) and \(\det(H)>0\), which implies the Hessian is negative definite. Therefore, \((x^\ast,y^\ast) = \big(\frac{110}{47}, \frac{60}{47}\big)\) is a strict local maximum. Since \(f\) is a quadratic with a negative definite Hessian everywhere, this is also the unique global maximum.

\bigskip

\subsection*{Problem 4 (16 points)}

\[
U(x,y)=x^{1/2}y^{1/2},\qquad p_x x + p_y y = I.
\]

\textbf{(a) Lagrangian and FOCs. (5 pts)}

Lagrangian:
\[
\mathcal{L}(x,y,\lambda) = x^{1/2}y^{1/2} - \lambda(p_x x + p_y y - I).
\]
First-order conditions:
\[
\frac{\partial \mathcal{L}}{\partial x}
= \frac{1}{2}x^{-1/2}y^{1/2} - \lambda p_x = 0,
\]
\[
\frac{\partial \mathcal{L}}{\partial y}
= \frac{1}{2}x^{1/2}y^{-1/2} - \lambda p_y = 0,
\]
\[
\frac{\partial \mathcal{L}}{\partial \lambda}
= -(p_x x + p_y y - I) = 0
\quad\Rightarrow\quad
p_x x + p_y y = I.
\]

\textbf{(b) Solve for \(x^\ast\) and \(y^\ast\). (5 pts)}

Divide the first FOC by the second (eliminate \(\lambda\)):
\[
\frac{\frac{1}{2}x^{-1/2}y^{1/2}}{\frac{1}{2}x^{1/2}y^{-1/2}}
= \frac{\lambda p_x}{\lambda p_y}
\quad\Rightarrow\quad
\frac{y/x}{1} = \frac{p_x}{p_y}.
\]
On the left:
\[
\frac{x^{-1/2}y^{1/2}}{x^{1/2}y^{-1/2}} 
= x^{-1} y = \frac{y}{x}.
\]
Hence:
\[
\frac{y}{x} = \frac{p_x}{p_y}
\quad\Rightarrow\quad
y = \frac{p_x}{p_y} x.
\]
Use the budget constraint:
\[
p_x x + p_y y = I
\quad\Rightarrow\quad
p_x x + p_y \left(\frac{p_x}{p_y}x\right) = p_x x + p_x x = 2 p_x x = I,
\]
so
\[
x^\ast = \frac{I}{2p_x}.
\]
Then
\[
y^\ast = \frac{p_x}{p_y} x^\ast
= \frac{p_x}{p_y}\cdot \frac{I}{2p_x}
= \frac{I}{2p_y}.
\]

\textbf{(c) Use the envelope theorem to compute \(\partial V / \partial I\). (3 pts)}

The indirect utility function is
\[
V(p_x,p_y,I) = U(x^\ast,y^\ast) 
= \left(\frac{I}{2p_x}\right)^{1/2}\left(\frac{I}{2p_y}\right)^{1/2}
= \frac{I}{2\sqrt{p_x p_y}}.
\]

By the envelope theorem for this equality-constrained maximization,
\[
\frac{\partial V}{\partial I} = \lambda^\ast,
\]
where \(\lambda^\ast\) is the optimal Lagrange multiplier.

We can compute directly from \(V\):
\[
\frac{\partial V}{\partial I} 
= \frac{1}{2\sqrt{p_x p_y}}.
\]
Thus \(\lambda^\ast = \dfrac{1}{2\sqrt{p_x p_y}}\).

\textbf{(d) Interpretation. (3 pts)}

\(\partial V/\partial I = \lambda^\ast\) is the \emph{marginal utility of income}: the increase in maximum attainable utility from a one-unit increase in income. It is positive and decreasing in the geometric mean of prices \(\sqrt{p_x p_y}\): higher prices reduce the marginal value of additional income.

\bigskip

\subsection*{Problem 5 (16 points)}

Capital evolves as
\[
\frac{dK}{dt} = I(t) - \delta K,\qquad I(t)=50 e^{0.01 t},\quad \delta = 0.05.
\]

\textbf{(a) Solve the differential equation for \(K(t)\). (6 pts)}

Rewrite in standard linear form:
\[
\frac{dK}{dt} + 0.05 K = 50 e^{0.01 t}.
\]
Integrating factor:
\[
\mu(t) = e^{\int 0.05 dt} = e^{0.05 t}.
\]
Multiply both sides by \(\mu(t)\):
\[
e^{0.05t}\frac{dK}{dt} + 0.05 e^{0.05 t} K = 50 e^{0.05 t} e^{0.01 t}
= 50 e^{0.06 t}.
\]
Left-hand side is derivative of \(K e^{0.05 t}\):
\[
\frac{d}{dt}(K e^{0.05 t}) = 50 e^{0.06 t}.
\]
Integrate:
\[
K(t)e^{0.05 t} = \int 50 e^{0.06 t} dt = \frac{50}{0.06} e^{0.06 t} + C
= \frac{2500}{3} e^{0.06 t} + C.
\]
Thus
\[
K(t) = \frac{2500}{3} e^{0.06 t} e^{-0.05 t} + C e^{-0.05 t}
= \frac{2500}{3} e^{0.01 t} + C e^{-0.05 t}.
\]

\textbf{(b) Particular solution with \(K(0)=100\). (4 pts)}

At \(t=0\),
\[
K(0) = \frac{2500}{3}\cdot 1 + C\cdot 1 = \frac{2500}{3} + C = 100.
\]
Hence
\[
C = 100 - \frac{2500}{3} = \frac{300 - 2500}{3} = -\frac{2200}{3}.
\]
So
\[
K(t) = \frac{2500}{3} e^{0.01 t} - \frac{2200}{3} e^{-0.05 t}.
\]

\textbf{(c) Compute \(\displaystyle \int_0^{20} I(t)\, dt\). (3 pts)}

\[
\int_0^{20} I(t)\, dt = \int_0^{20} 50 e^{0.01 t}\, dt
= 50 \int_0^{20} e^{0.01 t}\, dt
= 50 \left[\frac{1}{0.01} e^{0.01 t}\right]_0^{20}
= 50 \cdot 100 (e^{0.2} - 1).
\]
So
\[
\int_0^{20} I(t)\, dt = 5000 (e^{0.2} - 1).
\]

\textbf{(d) Long-run behavior of \(K(t)\). (3 pts)}

From
\[
K(t) = \frac{2500}{3} e^{0.01 t} - \frac{2200}{3} e^{-0.05 t},
\]
as \(t \to \infty\),
\[
e^{0.01 t} \to \infty,\quad e^{-0.05 t} \to 0.
\]
Thus
\[
K(t) \sim \frac{2500}{3} e^{0.01 t} \to \infty.
\]
Economically: because investment grows over time at a positive rate (1\%) and dominates depreciation in the long run, the capital stock keeps increasing without bound, growing asymptotically at net rate \(1\%\).

\end{document}
